\section{Personal Scenes}
\label{sec:personalScenes}

Start this chapter with utilizing the the Spiritual Realm to run several scenes tailored for each of the PCs. For each PC, there should be at least two scenes which put that character into a spotlight, allowing them to shine.

\subsubsection{Limited Rest}

After each scene, allow the PCs to take a Short Rest. To take a Long Rest, the party would need to enter a vision specifically designed for that purpose.

Unless there is a very good reason, Long Rest scenes should refrain from putting any PC into the spotlight. Save such opportunities for more meaningful scenes.

To make resource management valid in this part of the adventure, do not introduce Long Rest scenes too often - at most in between four other scenes.

\subsubsection{Hateful Ending}

At the end of each scene allow \npcMishram{}'s visage to appear in the vision. It may form from the clouds, rocks, shadows, or similar, and it will curse and make threats to the PCs.

Unless the PCs flee the vision soon after, they must make \skillCheck{20}{\skillDiscipline{}}. On a failure, they lose 1 Focus and 1 Maximum Focus until the next Long Rest.



\subsection{Endeavor Variant}

If your party likes to have a clear goal and tangible progress visualization, you may run this part as an Endeavor. To succeed on the Endeavor, the party must gather a number of successes equal to twice the number of PCs in the party.

Each Spiritual Realm encounter scene successfully resolved by the party awards them with one success towards the Endeavor progression.

There is no failing threshold for the Endeavor. If the PCs play enough encounters, they will eventually complete it. However, the longer they take, the more resources they will consume.





\subsection{Lead the Path}

Once a scene is over, an Enlightened Radiant may spend a Focus to make \skillCheck{25}{\skillLore{}}. On a success, the Radiant chooses the next scene the party encounters. That might be a specific event they experienced or a quiet moment of respite and rest.



\subsection{Example Scenes}

You should design the scenes for this part of the chapter to best align with the party, their story, and preferences of the players. It is up to you what scenes will be played. You can find some suggestions and examples in following sections:

\begin{itemize}
	\item \fullref{sec:fixTheFailure}
	\item \fullref{sec:backstoryFlashback}
	\item \fullref{sec:momentOfRespite}
	\item \fullref{sec:epicFight}
	\item \fullref{sec:unmadeProgression}
	\item \fullref{sec:familiarPast}
	\item \fullref{sec:learnTheLore}
	\item \fullref{sec:completeTheirStory}
\end{itemize}










\subsection{Aftermath}

Once you decide there is no need for more of the scenes to put the PCs in the spotlight, advance the party to level 15. Then, proceed to \fullref{sec:trailOfTheAges}.